\documentclass[12pt,a4paper]{article}
\usepackage[margin=2.5cm]{geometry}
\usepackage{amsmath,amssymb,amsfonts}
\usepackage{physics}
\usepackage{enumitem}
\usepackage{fancyhdr}
\usepackage{lastpage}
\usepackage{graphicx}
\usepackage{multicol}
\usepackage{array}
\usepackage{tabularx}
\usepackage{tikz}

% Page style
\pagestyle{fancy}
\fancyhf{}
\renewcommand{\headrulewidth}{0.4pt}
\renewcommand{\footrulewidth}{0.4pt}
\lhead{PHYS10012}
\rhead{Mock Examination}
\cfoot{Page \thepage\ of \pageref{LastPage}}

% Question numbering
\newcounter{questioncounter}
\newcommand{\question}[1]{%
  \stepcounter{questioncounter}%
  \vspace{0.5cm}%
  \noindent\textbf{Question \thequestioncounter.} \hfill [#1 marks]%
  \vspace{0.2cm}%
  \par%
}

% Sub-part environments
\newlist{parts}{enumerate}{2}
\setlist[parts,1]{label=(\alph*),leftmargin=2em,itemsep=0.3cm}
\setlist[parts,2]{label=(\roman*),leftmargin=2em,itemsep=0.2cm}

% Mark allocation in sub-parts
\renewcommand{\marks}[1]{\hfill [#1]}

% MCQ environment
\newcounter{mcqcounter}
\newcommand{\mcq}[1]{%
  \stepcounter{mcqcounter}%
  \vspace{0.3cm}%
  \noindent\textbf{\themcqcounter.} #1 \hfill [2 marks]%
  \vspace{0.1cm}%
}
\newlist{choices}{enumerate}{1}
\setlist[choices]{label=\Alph*.,leftmargin=2.5em,itemsep=0.1cm}

% Dirac notation shortcuts
\renewcommand{\ket}[1]{\left|#1\right\rangle}
\renewcommand{\bra}[1]{\left\langle #1\right|}
\renewcommand{\braket}[2]{\left\langle #1\middle|#2\right\rangle}
\renewcommand{\ketbra}[2]{\left|#1\right\rangle\!\left\langle #2\right|}
\renewcommand{\expect}[1]{\left\langle #1\right\rangle}

\begin{document}

% Title block
\begin{center}
{\Large\textbf{UNIVERSITY OF BRISTOL}}\\[0.3cm]
{\large Faculty of Science}\\[0.5cm]
{\Large\textbf{MOCK EXAMINATION --- Paper 1}}\\[0.3cm]
{\large\textbf{Core Physics I --- Classical, Quantum and Thermal Physics}}\\[0.2cm]
{\large Module Code: PHYS10012}\\[0.5cm]
{\large Duration: 3 Hours}\\[0.3cm]
{\large Total Marks: 100}\\[0.5cm]
\end{center}

\noindent\rule{\textwidth}{0.4pt}

\vspace{0.3cm}
\noindent\textbf{Instructions:}
\begin{itemize}[itemsep=0.1cm]
\item Answer \textbf{ALL} questions.
\item \textbf{Section A} (Oscillations and Waves): 10 multiple-choice questions, 20 marks total.
\item \textbf{Section B} (Quantum Physics): 10 multiple-choice questions, 20 marks total.
\item \textbf{Section C}: 4 long-answer questions, 60 marks total.
\item An approved calculator is permitted.
\item A formula sheet is provided at the end of this paper.
\item Show all working for Section C. Marks will be awarded for clear, logical solutions.
\end{itemize}

\noindent\rule{\textwidth}{0.4pt}
\newpage

% ============================================================
% SECTION A: OSCILLATIONS AND WAVES MCQ
% ============================================================
\section*{Section A: Oscillations and Waves}
\noindent\textit{Answer all 10 questions. Each question is worth 2 marks.}\\[0.3cm]

\setcounter{mcqcounter}{0}

% QUESTION A1: SHM — extract period/frequency from x(t)
\mcq{A particle undergoes simple harmonic motion described by $x(t) = 0.05\cos(6\pi t + \pi/4)$, where $x$ is in metres and $t$ in seconds. What is the period of oscillation?}
\begin{choices}
\item $\dfrac{1}{6}$ s
\item $\dfrac{1}{3}$ s
\item $3$ s
\item $6\pi$ s
\end{choices}

% QUESTION A2: SHM energy — find max velocity from given k and A
\mcq{A block of mass $0.4$ kg is attached to a spring with spring constant $k = 160$ N/m and oscillates with amplitude $A = 0.08$ m. What is the maximum speed of the block?}
\begin{choices}
\item $0.80$ m/s
\item $1.60$ m/s
\item $3.20$ m/s
\item $6.40$ m/s
\end{choices}

% QUESTION A3: Damped oscillation — identify damping regime
\mcq{A damped oscillator has mass $m = 0.25$ kg, spring constant $k = 100$ N/m, and damping coefficient $b = 10$ kg/s. The damping parameter is $\gamma = b/m$ and the natural frequency is $\omega_0 = \sqrt{k/m}$. What is the damping regime of this system?}
\begin{choices}
\item Underdamped ($\gamma < 2\omega_0$)
\item Critically damped ($\gamma = 2\omega_0$)
\item Overdamped ($\gamma > 2\omega_0$)
\item Undamped ($\gamma = 0$)
\end{choices}

% QUESTION A4: Forced oscillation — what happens at resonance (qualitative)
\mcq{A lightly damped oscillator is driven by an external sinusoidal force. At resonance, which of the following statements is correct?}
\begin{choices}
\item The amplitude is at its maximum and the phase lag of the displacement behind the driving force is $\pi/2$.
\item The amplitude is at its maximum and the displacement is in phase with the driving force.
\item The amplitude is zero and the phase lag is $\pi$.
\item The amplitude is at its maximum and the phase lag is $\pi$.
\end{choices}

% QUESTION A5: Wave speed on a string
\mcq{A guitar string has a linear mass density of $\mu = 0.004$ kg/m and is held under a tension of $T = 100$ N. What is the speed of transverse waves on this string?}
\begin{choices}
\item $25$ m/s
\item $50$ m/s
\item $158$ m/s
\item $250$ m/s
\end{choices}

% QUESTION A6: Impedance matching — when is there no reflection?
\mcq{Two strings are joined together. The impedance of string 1 is $Z_1$ and the impedance of string 2 is $Z_2$. Under what condition is there \textbf{no} reflected wave when a wave travelling on string 1 reaches the junction?}
\begin{choices}
\item $Z_1 = 0$
\item $Z_2 = 0$
\item $Z_1 = Z_2$
\item $Z_1 Z_2 = 1$
\end{choices}

% QUESTION A7: Sound intensity in decibels
\mcq{A jet engine at a certain distance produces a sound intensity of $I = 0.01$ W/m$^2$. What is the corresponding sound level in decibels? (Take $I_0 = 10^{-12}$ W/m$^2$.)}
\begin{choices}
\item $80$ dB
\item $90$ dB
\item $100$ dB
\item $120$ dB
\end{choices}

% QUESTION A8: Doppler effect — ambulance approaching
\mcq{An ambulance siren emits sound at a frequency of $f = 600$ Hz. The ambulance is approaching a stationary observer at a speed of $v_s = 34$ m/s. Taking the speed of sound in air as $v = 340$ m/s, what frequency does the observer hear?}
\begin{choices}
\item $540$ Hz
\item $600$ Hz
\item $660$ Hz
\item $667$ Hz
\end{choices}

% QUESTION A9: Standing waves on a string — fundamental frequency
\mcq{A string of length $L = 0.80$ m is fixed at both ends and supports transverse waves with speed $v = 320$ m/s. What is the fundamental frequency of vibration?}
\begin{choices}
\item $100$ Hz
\item $200$ Hz
\item $400$ Hz
\item $640$ Hz
\end{choices}

% QUESTION A10: Beats
\mcq{Two tuning forks are sounded together. Fork A has a frequency of $440$ Hz and fork B has a frequency of $446$ Hz. How many beats per second are heard?}
\begin{choices}
\item $3$
\item $6$
\item $443$
\item $886$
\end{choices}

\newpage

% ============================================================
% SECTION B: QUANTUM PHYSICS MCQ
% ============================================================
\section*{Section B: Quantum Physics}
\noindent\textit{Answer all 10 questions. Each question is worth 2 marks.}\\[0.3cm]

\setcounter{mcqcounter}{0}

% QUESTION B1: Identify which expression is a valid ket/bra/operator
\mcq{Which of the following expressions represents a \textbf{linear operator} acting on a Hilbert space?}
\begin{choices}
\item $\braket{\alpha}{\beta}$
\item $\ketbra{\alpha}{\beta}$
\item $\ket{\alpha}$
\item $\bra{\beta}$
\end{choices}

% QUESTION B2: Normalisation — find normalisation constant
\mcq{A quantum state is written as $\ket{\psi} = \alpha\bigl(\ket{1} + 2\ket{2}\bigr)$, where $\ket{1}$ and $\ket{2}$ are orthonormal basis states and $\alpha$ is a real, positive normalisation constant. What is the value of $\alpha$?}
\begin{choices}
\item $1/3$
\item $1/\sqrt{3}$
\item $1/\sqrt{5}$
\item $1/5$
\end{choices}

% QUESTION B3: Inner product calculation
\mcq{Let $\ket{\phi} = \dfrac{1}{\sqrt{2}}\ket{1} + \dfrac{i}{\sqrt{2}}\ket{2}$ and $\ket{\psi} = \dfrac{1}{\sqrt{2}}\ket{1} - \dfrac{i}{\sqrt{2}}\ket{2}$, where $\{\ket{1}, \ket{2}\}$ is an orthonormal basis. What is $\braket{\phi}{\psi}$?}
\begin{choices}
\item $0$
\item $1$
\item $i$
\item $-i$
\end{choices}

% QUESTION B4: Hermitian operator test
\mcq{Consider the following operators defined on a two-state orthonormal basis $\{\ket{1}, \ket{2}\}$. Which one is Hermitian?}
\begin{choices}
\item $\hat{A} = i\ketbra{1}{2} + i\ketbra{2}{1}$
\item $\hat{B} = \ketbra{1}{2} - \ketbra{2}{1}$
\item $\hat{C} = \ketbra{1}{2} + \ketbra{2}{1}$
\item $\hat{D} = i\ketbra{1}{2} - i\ketbra{2}{1}$
\end{choices}

% QUESTION B5: Born rule — given state in energy eigenbasis, find P(E1)
\mcq{A particle is in the state $\ket{\psi} = \dfrac{3}{5}\ket{E_1} + \dfrac{4}{5}\ket{E_2}$, where $\ket{E_1}$ and $\ket{E_2}$ are orthonormal energy eigenstates with energies $E_1$ and $E_2$ respectively. If the energy is measured, what is the probability of obtaining $E_1$?}
\begin{choices}
\item $3/5$
\item $9/25$
\item $16/25$
\item $4/5$
\end{choices}

% QUESTION B6: Spin measurement — particle in |↑z⟩, measure S_x
\mcq{A spin-1/2 particle is prepared in the state $\ket{\!\uparrow_z}$. A measurement of $S_x$ is performed. What is the probability of obtaining $+\hbar/2$?}
\begin{choices}
\item $0$
\item $1/4$
\item $1/2$
\item $1$
\end{choices}

% QUESTION B7: Time evolution — which quantity is conserved
\mcq{A quantum system has Hamiltonian $H$. A state is prepared at $t=0$ and then evolves in time under $H$. Which of the following quantities is \textbf{always} conserved during this time evolution?}
\begin{choices}
\item The probabilities of measuring each position eigenvalue
\item The expectation value of $S_x$
\item The probability of measuring each energy eigenvalue
\item The expectation value of $S_z$
\end{choices}

% QUESTION B8: Entanglement — identify which 2-particle state is entangled
\mcq{Consider two spin-1/2 particles, labelled 1 and 2. Which of the following two-particle states is \textbf{entangled}?}
\begin{choices}
\item $\ket{\!\uparrow_z}_1 \ket{\!\downarrow_z}_2$
\item $\dfrac{1}{\sqrt{2}}\bigl(\ket{\!\uparrow_z}_1 + \ket{\!\downarrow_z}_1\bigr)\ket{\!\uparrow_z}_2$
\item $\dfrac{1}{\sqrt{2}}\bigl(\ket{\!\uparrow_z}_1\ket{\!\downarrow_z}_2 - \ket{\!\downarrow_z}_1\ket{\!\uparrow_z}_2\bigr)$
\item $\ket{\!\downarrow_z}_1 \ket{\!\downarrow_z}_2$
\end{choices}

% QUESTION B9: Identical particles — swapping two fermions
\mcq{Two identical fermions occupy a joint quantum state $\ket{\Psi}$. What happens to $\ket{\Psi}$ when the labels of the two particles are exchanged (i.e.\ the SWAP operator is applied)?}
\begin{choices}
\item $\ket{\Psi}$ is unchanged: $\text{SWAP}\ket{\Psi} = +\ket{\Psi}$
\item $\ket{\Psi}$ acquires a factor of $-1$: $\text{SWAP}\ket{\Psi} = -\ket{\Psi}$
\item $\ket{\Psi}$ acquires a factor of $i$: $\text{SWAP}\ket{\Psi} = i\ket{\Psi}$
\item $\ket{\Psi}$ becomes orthogonal to the original state
\end{choices}

% QUESTION B10: Particle physics — quark content of neutron
\mcq{A neutron is a composite particle made of quarks. What is the quark content of a neutron?}
\begin{choices}
\item $uuu$
\item $uud$
\item $udd$
\item $ddd$
\end{choices}

\newpage

% ============================================================
% SECTION C: LONG ANSWER QUESTIONS
% ============================================================
\section*{Section C: Long Answer Questions}
\noindent\textit{Answer all 4 questions. Show all working.}

\setcounter{questioncounter}{0}

% ---- QUESTION C1: MECHANICS (15 marks) ----
\question{15}

\noindent\textbf{Projectile Motion and Energy Conservation}

\noindent A ball is launched from ground level at speed $v_0 = 25$ m/s at an angle $\theta = 53^\circ$ above the horizontal. Take $g = 9.81$ m/s$^2$.

\begin{parts}
\item Find the maximum height $H$ reached by the ball. \qmarks{4}

\item Find the range $R$ of the ball (the horizontal distance when it returns to ground level). \qmarks{3}

\item A second ball of mass $m = 0.2$ kg is dropped from rest at height $H$ (the value found in part~(a)). It falls onto a vertical spring of spring constant $k = 500$ N/m resting on the ground. Using conservation of energy, find the maximum compression $x$ of the spring. \qmarks{4}

\item The fully compressed spring is now used to launch a marble of mass $m_{\text{marble}} = 0.1$ kg horizontally from the top of a $20$ m high cliff. Using the compression found in part~(c), find the speed of the marble at the instant it hits the ground. \qmarks{4}
\end{parts}

\newpage

% ---- QUESTION C2: PROPERTIES OF MATTER (15 marks) ----
\question{15}

\noindent\textbf{Thermodynamic Processes}

\noindent Two moles ($n = 2$) of a monatomic ideal gas are initially at temperature $T_1 = 300$ K and volume $V_1 = 0.05$ m$^3$. Take $R = 8.314$ J/(mol$\cdot$K), $\gamma = 5/3$ for a monatomic ideal gas, and $C_V = \frac{3}{2}R$ per mole.

\begin{parts}
\item The gas undergoes a \textbf{reversible isothermal expansion} to volume $V_2 = 0.15$ m$^3$. Calculate the work done by the gas. \qmarks{3}

\item Calculate the heat absorbed by the gas during this isothermal expansion. \qmarks{3}

\item The gas then undergoes a \textbf{reversible adiabatic expansion} from $(T_2 = 300\text{ K},\; V_2 = 0.15\text{ m}^3)$ to a final volume $V_3 = 0.30$ m$^3$. Find the final temperature $T_3$. \qmarks{4}

\item Calculate the work done by the gas during the adiabatic expansion. \qmarks{2}

\item Calculate the total entropy change of the universe for the isothermal process in part~(a). Comment on whether this result is consistent with the process being reversible. \qmarks{3}
\end{parts}

\newpage

% ---- QUESTION C3: OSCILLATIONS AND WAVES (15 marks) ----
\question{15}

\noindent\textbf{Damped and Forced Oscillations}

\begin{parts}
\item Write down the equation of motion for a damped harmonic oscillator of mass $m$, spring constant $k$, and damping coefficient $b$ (where the damping force is $-b\dot{x}$). Define the damping parameter $\gamma$ in terms of $b$ and $m$. \qmarks{2}

\item A block of mass $m = 0.5$ kg is attached to a spring with spring constant $k = 200$ N/m and experiences a damping coefficient $b = 2$ kg/s. Calculate the natural angular frequency $\omega_0$, the damping parameter $\gamma$, and the damped angular frequency $\omega_d$. \qmarks{3}

\item Calculate the quality factor $Q$ of this system. State whether the system is underdamped, critically damped, or overdamped, giving a brief justification. \qmarks{2}

\item The system is now driven by an external force $F(t) = F_0\cos(\omega t)$ with $F_0 = 5$ N. Find the steady-state amplitude of oscillation when the driving angular frequency is $\omega = 19$ rad/s. \qmarks{4}

\item Find the driving angular frequency $\omega_{\text{res}}$ at which the steady-state amplitude is maximum (the amplitude resonance frequency). Calculate the maximum amplitude $A_{\text{max}}$ at this resonance frequency. \qmarks{4}
\end{parts}

\newpage

% ---- QUESTION C4: QUANTUM PHYSICS (15 marks) ----
\question{15}

\noindent\textbf{Spin-1/2 and Measurement}

\noindent A spin-1/2 particle is prepared in the state
\[
\ket{\psi} = \frac{1}{\sqrt{3}}\ket{\!\uparrow_z} + \sqrt{\frac{2}{3}}\,\ket{\!\downarrow_z}.
\]

\begin{parts}
\item Verify that $\ket{\psi}$ is normalised. If $S_z$ is measured, state the possible outcomes and their probabilities. \qmarks{3}

\item Express $\ket{\psi}$ in the $S_x$ eigenbasis $\{\ket{\!\uparrow_x},\, \ket{\!\downarrow_x}\}$. Hence, if $S_x$ is measured on the original state $\ket{\psi}$, find the probability of each outcome. \qmarks{4}

\item After measuring $S_x$ and obtaining the result $+\hbar/2$, the particle's state collapses to $\ket{\!\uparrow_x}$. If $S_z$ is then measured on this post-measurement state, find the probability of each outcome. \qmarks{4}

\item Consider the operator $\hat{A} = \hbar\bigl(\ketbra{\!\uparrow_z}{\downarrow_z\!} + \ketbra{\!\downarrow_z}{\uparrow_z\!}\bigr)$. Find the eigenvalues and normalised eigenvectors of $\hat{A}$. Determine whether $\hat{A}$ is Hermitian, justifying your answer. \qmarks{4}
\end{parts}

\newpage

% ============================================================
% FORMULA SHEET
% ============================================================
% EQUATION SHEET — PHYS10012 Core Physics I
% This file is \input{} at the end of each exam paper

\newpage
\section*{Formula Sheet}
\noindent\rule{\textwidth}{0.4pt}

\subsection*{Mathematical Formulae}

\begin{align*}
&\text{Binomial expansion:} & (1+x)^n &\approx 1 + nx + \frac{n(n-1)}{2!}x^2 + \cdots \quad (|x| \ll 1) \\[4pt]
&\text{Euler's formula:} & e^{i\theta} &= \cos\theta + i\sin\theta \\[4pt]
&\text{Trig identities:} & \sin A + \sin B &= 2\cos\!\left(\tfrac{A-B}{2}\right)\sin\!\left(\tfrac{A+B}{2}\right) \\
& & \cos A + \cos B &= 2\cos\!\left(\tfrac{A-B}{2}\right)\cos\!\left(\tfrac{A+B}{2}\right)
\end{align*}

\subsection*{Mechanics}

\begin{align*}
&\text{SUVAT:} & v &= u + at, \quad s = ut + \tfrac{1}{2}at^2, \quad v^2 = u^2 + 2as, \quad s = \tfrac{1}{2}(u+v)t \\[4pt]
&\text{Dot product:} & \mathbf{a}\cdot\mathbf{b} &= |\mathbf{a}||\mathbf{b}|\cos\theta = a_x b_x + a_y b_y + a_z b_z \\[4pt]
&\text{Cross product:} & \mathbf{a}\times\mathbf{b} &= |\mathbf{a}||\mathbf{b}|\sin\theta\;\hat{\mathbf{n}} = \begin{vmatrix}\hat{\mathbf{i}}&\hat{\mathbf{j}}&\hat{\mathbf{k}}\\a_x&a_y&a_z\\b_x&b_y&b_z\end{vmatrix} \\[4pt]
&\text{Rocket equation:} & \Delta v &= u \ln\!\left(\frac{m_0}{m_f}\right) \\[4pt]
&\text{Gravitational PE:} & U &= -\frac{GMm}{r} \\[4pt]
&\text{Escape velocity:} & v_e &= \sqrt{\frac{2GM}{R}} \\[4pt]
&\text{Force from PE:} & F &= -\frac{dU}{dx} \\[4pt]
&\text{Elastic collision (1D):} & v_1' &= \frac{m_1 - m_2}{m_1 + m_2}v_1 + \frac{2m_2}{m_1+m_2}v_2 \\[4pt]
&\text{Centre of mass:} & \mathbf{R}_\text{cm} &= \frac{\sum m_i\mathbf{r}_i}{\sum m_i} \\[4pt]
&\text{Reduced mass:} & \mu &= \frac{m_1 m_2}{m_1 + m_2} \\[4pt]
&\text{Centripetal accel:} & a_c &= \frac{v^2}{r} = \omega^2 r \\[4pt]
&\text{Parallel axis thm:} & I &= I_\text{cm} + Md^2 \\[4pt]
&\text{Atwood machine:} & a &= \frac{(m_1-m_2)g}{m_1+m_2}, \quad T = \frac{2m_1 m_2 g}{m_1+m_2}
\end{align*}

\noindent\textbf{Moments of inertia} (about axis through centre of mass):
\begin{center}
\begin{tabular}{ll}
Thin ring (radius $R$): & $I = MR^2$ \\
Solid disk (radius $R$): & $I = \tfrac{1}{2}MR^2$ \\
Solid sphere (radius $R$): & $I = \tfrac{2}{5}MR^2$ \\
Hollow sphere (radius $R$): & $I = \tfrac{2}{3}MR^2$ \\
Thin rod (length $L$, about centre): & $I = \tfrac{1}{12}ML^2$ \\
Thin rod (length $L$, about end): & $I = \tfrac{1}{3}ML^2$ \\
\end{tabular}
\end{center}

\subsection*{Properties of Matter}

\begin{align*}
&\text{Ideal gas:} & PV &= nRT \\[4pt]
&\text{Heat capacities:} & C_P - C_V &= R \quad\text{(per mole)} \\[4pt]
&\text{Adiabatic:} & PV^\gamma &= \text{const}, \quad TV^{\gamma-1} = \text{const} \\[4pt]
&\text{Isothermal work:} & W &= nRT\ln\!\left(\frac{V_2}{V_1}\right) \\[4pt]
&\text{Carnot efficiency:} & \eta &= 1 - \frac{T_c}{T_h} \\[4pt]
&\text{Entropy:} & dS &= \frac{dQ_\text{rev}}{T}, \quad \Delta S = mc\ln\!\left(\frac{T_2}{T_1}\right), \quad \Delta S = nR\ln\!\left(\frac{V_2}{V_1}\right) \\[4pt]
&\text{Lennard-Jones:} & V(r) &= 4\varepsilon\left[\left(\frac{a}{r}\right)^{12} - \left(\frac{a}{r}\right)^{6}\right] \\[4pt]
&\text{Mean KE:} & \langle KE \rangle &= \tfrac{3}{2}k_B T \\[4pt]
&\text{MB speeds:} & c_\text{mp} &= \sqrt{\frac{2k_BT}{m}}, \quad c_\text{avg} = \sqrt{\frac{8k_BT}{\pi m}}, \quad c_\text{rms} = \sqrt{\frac{3k_BT}{m}}
\end{align*}

\subsection*{Oscillations and Waves}

\begin{align*}
&\text{SHM:} & x(t) &= A\cos(\omega_0 t + \phi), \quad \omega_0 = \sqrt{k/m} \\[4pt]
&\text{Damped SHM:} & x(t) &= Ae^{-\gamma t/2}\cos(\omega_d t + \phi), \quad \omega_d = \sqrt{\omega_0^2 - \gamma^2/4} \\[4pt]
&\text{Quality factor:} & Q &= \omega_0/\gamma \\[4pt]
&\text{Forced amplitude:} & A(\omega) &= \frac{F_0/m}{\sqrt{(\omega_0^2 - \omega^2)^2 + \gamma^2\omega^2}} \\[4pt]
&\text{Phase lag:} & \tan\delta &= \frac{\gamma\omega}{\omega_0^2 - \omega^2} \\[4pt]
&\text{Pendulums:} & T_\text{simple} &= 2\pi\sqrt{L/g}, \quad T_\text{physical} = 2\pi\sqrt{I/(Mgh)} \\[4pt]
&\text{Wave equation:} & \frac{\partial^2 y}{\partial x^2} &= \frac{1}{v^2}\frac{\partial^2 y}{\partial t^2} \\[4pt]
&\text{String speed:} & v &= \sqrt{T/\mu} \\[4pt]
&\text{Sound in gas:} & v &= \sqrt{\gamma P/\rho} \\[4pt]
&\text{Harmonic wave:} & y &= A\sin(kx - \omega t) \\[4pt]
&\text{String impedance:} & Z &= \mu v = \sqrt{\mu T} \\[4pt]
&\text{Reflection:} & r &= \frac{Z_1 - Z_2}{Z_1 + Z_2}, \quad t = \frac{2Z_1}{Z_1+Z_2} \\[4pt]
&\text{Power (string):} & \langle P\rangle &= \tfrac{1}{2}\mu\omega^2 A^2 v \\[4pt]
&\text{Decibels:} & \beta &= 10\log_{10}(I/I_0), \quad I_0 = 10^{-12}\;\text{W/m}^2 \\[4pt]
&\text{Doppler:} & f' &= f\frac{v \pm v_o}{v \mp v_s} \\[4pt]
&\text{Group velocity:} & v_g &= \frac{d\omega}{dk} \\[4pt]
&\text{String harmonics:} & f_n &= \frac{nv}{2L} \quad (n=1,2,3,\ldots) \\[4pt]
&\text{Open-closed pipe:} & f_n &= \frac{nv}{4L} \quad (n=1,3,5,\ldots)
\end{align*}

\subsection*{Quantum Physics}

\begin{align*}
&\text{Schr\"{o}dinger eq:} & i\hbar\frac{d}{dt}\ket{\psi(t)} &= H\ket{\psi(t)} \\[4pt]
&\text{Time evolution:} & \ket{\psi(t)} &= \sum_n c_n e^{-iE_n t/\hbar}\ket{E_n}, \quad U(t) = e^{-iHt/\hbar} \\[4pt]
&\text{Spin operators:} & S_z &= \frac{\hbar}{2}\bigl(\ketbra{\!\uparrow}{\uparrow\!} - \ketbra{\!\downarrow}{\downarrow\!}\bigr) \\
& & S_x &= \frac{\hbar}{2}\bigl(\ketbra{\!\uparrow}{\downarrow\!} + \ketbra{\!\downarrow}{\uparrow\!}\bigr) \\
& & S_y &= \frac{\hbar}{2}\bigl(-i\ketbra{\!\uparrow}{\downarrow\!} + i\ketbra{\!\downarrow}{\uparrow\!}\bigr) \\[4pt]
&\text{Spin eigenstates:} & \ket{\!\uparrow_x} &= \tfrac{1}{\sqrt{2}}\bigl(\ket{\!\uparrow_z} + \ket{\!\downarrow_z}\bigr), \quad \ket{\!\downarrow_x} = \tfrac{1}{\sqrt{2}}\bigl(\ket{\!\uparrow_z} - \ket{\!\downarrow_z}\bigr) \\
& & \ket{\!\uparrow_y} &= \tfrac{1}{\sqrt{2}}\bigl(\ket{\!\uparrow_z} + i\ket{\!\downarrow_z}\bigr), \quad \ket{\!\downarrow_y} = \tfrac{1}{\sqrt{2}}\bigl(\ket{\!\uparrow_z} - i\ket{\!\downarrow_z}\bigr) \\[4pt]
&\text{Pauli matrices:} & \sigma_x &= \begin{pmatrix}0&1\\1&0\end{pmatrix}, \quad \sigma_y = \begin{pmatrix}0&-i\\i&0\end{pmatrix}, \quad \sigma_z = \begin{pmatrix}1&0\\0&-1\end{pmatrix}
\end{align*}

\subsection*{Physical Constants}

\begin{center}
\begin{tabular}{lll}
Speed of light & $c$ & $3.00 \times 10^{8}$ m/s \\
Planck's constant & $h$ & $6.626 \times 10^{-34}$ J$\cdot$s \\
Reduced Planck's constant & $\hbar$ & $1.055 \times 10^{-34}$ J$\cdot$s \\
Boltzmann constant & $k_B$ & $1.381 \times 10^{-23}$ J/K \\
Gas constant & $R$ & $8.314$ J/(mol$\cdot$K) \\
Avogadro's number & $N_A$ & $6.022 \times 10^{23}$ mol$^{-1}$ \\
Gravitational constant & $G$ & $6.674 \times 10^{-11}$ N$\cdot$m$^2$/kg$^2$ \\
Acceleration due to gravity & $g$ & $9.81$ m/s$^2$ \\
Electron mass & $m_e$ & $9.109 \times 10^{-31}$ kg \\
Proton mass & $m_p$ & $1.673 \times 10^{-27}$ kg \\
Elementary charge & $e$ & $1.602 \times 10^{-19}$ C \\
Mass of Earth & $M_\oplus$ & $5.97 \times 10^{24}$ kg \\
Radius of Earth & $R_\oplus$ & $6.37 \times 10^{6}$ m \\
\end{tabular}
\end{center}

\noindent\rule{\textwidth}{0.4pt}


\end{document}
